\labsection{10}{Reinforcement learning and final integration}{sec:lab10-reinforcement}

\begin{labproject}{Learn a policy by interaction, then connect the result to the full course workflow}
\textbf{Coverage.} Chapter 20 with a handoff to the capstone.\\
\textbf{Core environment.} A self-contained $4\times4$ GridWorld.\\
\textbf{Goal.} Train a tabular Q-learning agent, evaluate its greedy policy separately from exploratory training, and articulate how reinforcement learning changes the data/target/evaluation assumptions used earlier in the course.

\begin{labmode}
This is the tenth and final scheduled lab. The capstone that follows is an integration project, not Lab 11.
\end{labmode}

\smallskip\noindent\textit{Implementation note.} Code is shown in notebook-sized blocks rather than as one complete listing. Run the blocks in order; each framed block corresponds to a canonical Jupyter code cell.\par

\medskip\noindent\textbf{Part A --- Tabular Q-learning in GridWorld.}\par
\lstinputlisting[style=mlpython]{jupyter_cells/lab10-reinforcement_integration/block01.py}
\lstinputlisting[style=mlpython]{jupyter_cells/lab10-reinforcement_integration/block02.py}
\lstinputlisting[style=mlpython]{jupyter_cells/lab10-reinforcement_integration/block03.py}
\lstinputlisting[style=mlpython]{jupyter_cells/lab10-reinforcement_integration/block04.py}
\lstinputlisting[style=mlpython]{jupyter_cells/lab10-reinforcement_integration/block05.py}
\lstinputlisting[style=mlpython]{jupyter_cells/lab10-reinforcement_integration/block06.py}
\lstinputlisting[style=mlpython]{jupyter_cells/lab10-reinforcement_integration/block07.py}

\begin{tryit}
Change one of reward, trap location, exploration schedule, or discount factor. Predict the behavioral consequence before retraining and compare the learned greedy policy afterward.
\end{tryit}
\end{labproject}

\begin{labdeliverable}
Submit the reward specification, one pre-training prediction, the learning curve, repeated greedy-policy success rate, and one controlled environment or hyperparameter change. End with a short capstone handoff naming the learning signal, validation boundary, baseline, and metric you would choose for a new problem.
\end{labdeliverable}
